%# -*- coding: utf-8-unix -*-

\chapter*{PREFACE}
\addcontentsline{toc}{chapter}{PREFACE}

\begin{quote}
\centering\itshape
Fermat, Euler, Lagrange, Legendre... introitum ad penetralia huius divinae scientiae aperuerunt, quantisque divitiis abundent patefecerunt. \\
\hfill --- Gauss, \textit{Disquisitiones Arithmeticae}
\end{quote}

\vspace{1em}

The study of transcendental numbers, springing from such diverse sources as the ancient Greek question concerning the squaring of the circle, the rudimentary researches of Liouville and Cantor, Hermite's investigations on the exponential function and the seventh of Hilbert's famous list of 23 problems, has now developed into a fertile and extensive theory, enriching widespread branches of mathematics; and the time has seemed opportune to prepare a systematic treatise. My aim has been to provide a comprehensive account of the recent major discoveries in the field; the text includes, more especially, expositions of the latest theories relating to linear forms in the logarithms of algebraic numbers, of Schmidt's generalization of the Thue-Siegel-Roth theorem, of Shidlovsky's work on Siegel's $E$-functions and of Sprindzuk's solution to the Mahler conjecture. Classical aspects of the subject are discussed in the course of the narrative; in particular, to facilitate the acquisition of a true historical perspective, a survey of the theory as it existed at about the turn of the century is given at the beginning. Proofs in the subject tend, as will be appreciated, to be long and intricate, and thus it has been necessary to select for detailed treatment only the most fundamental results; moreover, generally speaking, emphasis has been placed on arguments which have led to the strongest propositions known to date or have yielded the widest application. Nevertheless, it is hoped that adequate references have been included to associated works.

Notwithstanding its long history, it will be apparent that the theory of transcendental numbers bears a youthful countenance. Many topics would certainly benefit by deeper studies and several famous longstanding problems remain open. As examples, one need mention only the celebrated conjectures concerning the algebraic independence of $e$ and $\pi$ and the transcendence of Euler's constant $\gamma$, the solution to either of which would represent a major advance. If this book should play some small role in promoting future progress, the author will be well satisfied.

The text has arisen from numerous lectures delivered in Cambridge, America and elsewhere, and it has also formed the substance of an Adams Prize essay.

I am grateful to Dr D. W. Masser for his kind assistance in checking the proofs, and also to the Cambridge University Press for the care they have taken with the printing.

\vspace{1em}

\noindent
\textit{Cambridge, 1974} \hfill \textbf{A. B.}
